
\section{The Question Is Older Than It Looks}

God brings Abram outside and commands him to look toward heaven and count the
stars (Gen 15:5). Seven verses later, the sun is about to set. At verse 17 it
has set and darkness has come. If the chapter narrates one continuous episode
in sequence, the star command occurs before sunset.

The observation is often treated as a modern cleverness: perhaps the stars
were present but invisible, making the command a demand for faith. That
theological reading may be fruitful, but first the historical question must be
asked. Did earlier readers notice the apparent daylight? Did they preserve one
episode or insert a temporal break? Did ancient versions alter the clock? And
does the grammar force continuity?

Don Isaac Abarbanel proves that the difficulty is not new. His seventh question
before commenting on Genesis 15 states the dilemma: the event cannot be day
because stars are not visible by day, yet cannot be night because the text
later says the sun was about to set. Abarbanel's formulation matters because it
prevents the inquiry from being dismissed as a novelty. His answer is not the
only one, and this study does not recruit him for a predetermined conclusion.

